% math0 manuscript. Build with: latexmk -pdf paper.tex
\documentclass[11pt]{amsart}
\usepackage[margin=1.1in]{geometry}
\usepackage{amsmath,amssymb,amsthm,mathtools}
\usepackage{url}
\usepackage{hyperref}
\usepackage{xcolor}
\usepackage{listings}
\usepackage{newunicodechar}
\usepackage{booktabs}

\lstdefinelanguage{lean}{
  keywords={theorem,lemma,def,abbrev,example,fun,by,have,show,from,let,in,if,then,else,match,with,
    namespace,end,open,import,section,variable,instance,structure,where,calc,obtain,intro,exact,
    refine,rcases,rintro,simp,omega,decide,norm_num,linarith,nlinarith,positivity,induction,cases,
    private,unfold,rw,at},
  sensitive=true, comment=[l]{--}, morecomment=[s]{/-}{-/}, string=[b]"}
\lstset{language=lean, basicstyle=\ttfamily\footnotesize, columns=fullflexible, keepspaces=true, breaklines=true,
  keywordstyle=\color{blue!60!black}\bfseries, commentstyle=\color{gray}, frame=single,
  extendedchars=true, inputencoding=utf8,
  literate={ℕ}{{$\mathbb{N}$}}1 {∀}{{$\forall$}}1 {↔}{{$\leftrightarrow$}}1 {¬}{{$\neg$}}1
    {←}{{$\leftarrow$}}1 {≡}{{$\equiv$}}1 {∣}{{$\mid$}}1 {·}{{$\cdot$}}1}
% Unicode commonly found in Lean sources.
\newunicodechar{ℕ}{\ensuremath{\mathbb{N}}}
\newunicodechar{ℤ}{\ensuremath{\mathbb{Z}}}
\newunicodechar{ℚ}{\ensuremath{\mathbb{Q}}}
\newunicodechar{ℝ}{\ensuremath{\mathbb{R}}}
\newunicodechar{ℂ}{\ensuremath{\mathbb{C}}}
\newunicodechar{∀}{\ensuremath{\forall}}
\newunicodechar{∃}{\ensuremath{\exists}}
\newunicodechar{→}{\ensuremath{\to}}
\newunicodechar{←}{\ensuremath{\leftarrow}}
\newunicodechar{↔}{\ensuremath{\leftrightarrow}}
\newunicodechar{↦}{\ensuremath{\mapsto}}
\newunicodechar{¬}{\ensuremath{\neg}}
\newunicodechar{∧}{\ensuremath{\wedge}}
\newunicodechar{∨}{\ensuremath{\vee}}
\newunicodechar{≤}{\ensuremath{\le}}
\newunicodechar{≥}{\ensuremath{\ge}}
\newunicodechar{≠}{\ensuremath{\ne}}
\newunicodechar{∈}{\ensuremath{\in}}
\newunicodechar{∉}{\ensuremath{\notin}}
\newunicodechar{⊆}{\ensuremath{\subseteq}}
\newunicodechar{∩}{\ensuremath{\cap}}
\newunicodechar{∪}{\ensuremath{\cup}}
\newunicodechar{∑}{\ensuremath{\sum}}
\newunicodechar{∏}{\ensuremath{\prod}}
\newunicodechar{⟨}{\ensuremath{\langle}}
\newunicodechar{⟩}{\ensuremath{\rangle}}
\newunicodechar{·}{\ensuremath{\cdot}}
\newunicodechar{∘}{\ensuremath{\circ}}
\newunicodechar{⁻}{\ensuremath{^{-}}}
\newunicodechar{¹}{\ensuremath{^{1}}}
\newunicodechar{²}{\ensuremath{^{2}}}
\newunicodechar{₀}{\ensuremath{_{0}}}
\newunicodechar{₁}{\ensuremath{_{1}}}
\newunicodechar{₂}{\ensuremath{_{2}}}
\newunicodechar{α}{\ensuremath{\alpha}}
\newunicodechar{β}{\ensuremath{\beta}}
\newunicodechar{ε}{\ensuremath{\varepsilon}}
\newunicodechar{δ}{\ensuremath{\delta}}
\newunicodechar{π}{\ensuremath{\pi}}
\newunicodechar{σ}{\ensuremath{\sigma}}
\newunicodechar{φ}{\ensuremath{\varphi}}
\newunicodechar{∞}{\ensuremath{\infty}}
\newunicodechar{⌊}{\ensuremath{\lfloor}}
\newunicodechar{⌋}{\ensuremath{\rfloor}}
\newunicodechar{⌈}{\ensuremath{\lceil}}
\newunicodechar{⌉}{\ensuremath{\rceil}}
\newunicodechar{∣}{\ensuremath{\mid}}
\newunicodechar{≡}{\ensuremath{\equiv}}
\newunicodechar{×}{\ensuremath{\times}}
\newunicodechar{⧸}{\ensuremath{/}}
\newunicodechar{‖}{\ensuremath{\|}}
\newunicodechar{√}{\ensuremath{\surd}}
\newunicodechar{∅}{\ensuremath{\emptyset}}
\newunicodechar{⦃}{\ensuremath{\{\!|}}
\newunicodechar{⦄}{\ensuremath{|\!\}}}
\newunicodechar{𝓝}{\ensuremath{\mathcal{N}}}
\newunicodechar{≫}{\ensuremath{\gg}}
\newunicodechar{≪}{\ensuremath{\ll}}

\theoremstyle{plain}
\newtheorem{theorem}{Theorem}
\newtheorem{lemma}[theorem]{Lemma}
\newtheorem{proposition}[theorem]{Proposition}
\newtheorem{corollary}[theorem]{Corollary}
\newtheorem{conjecture}[theorem]{Conjecture}
\theoremstyle{definition}
\newtheorem{definition}[theorem]{Definition}
\newtheorem{remark}[theorem]{Remark}

\emergencystretch=3em
\raggedbottom
% Lean identifiers in running text (breakable at dots and underscores).
\expandafter\def\expandafter\UrlBreaks\expandafter{\UrlBreaks\do\_\do\.}
\DeclareUrlCommand\decl{\urlstyle{tt}}
\newcommand{\N}{\mathbb{N}}

\title[A counterexample to the squarefree conjecture in OEIS A083020]{The number $3^{3^{18}}+2$ is divisible by $31^2$: a counterexample to Hilliard's squarefree conjecture in OEIS A083020}
\author{math0}
\address{Produced by Claude (an AI model developed by Anthropic) in the math0 project}
\date{October 21, 2026}

\begin{document}
\begin{abstract}
The OEIS entry A083020 (Cino Hilliard, 2003) lists the prime factors of the numbers $N_n=3^{3^n}+2$, $n\ge0$, and conjectures that every $N_n$ is squarefree.
We show that $31^2$ divides $N_{18}$, a number with $184{,}846{,}550$ decimal digits, so the conjecture is false; the proof reduces the exponent $3^{18}$ modulo $\varphi(31^2)=930$ and then needs only the congruence $3^{159}\equiv-2\pmod{961}$.
We also show, without formalizing it, that $31^2\mid N_n$ exactly when $n\equiv18\pmod{60}$ and $29^2\mid N_n$ exactly when $n\equiv73\pmod{84}$, so infinitely many $N_n$ are not squarefree.
To our knowledge the counterexample is new: the OEIS entry still states the conjecture and we found no recorded resolution, although the argument is elementary.
The refutation is verified in Lean~4 with Mathlib against the frozen statement that math0 formalized from the OEIS entry; its definitions are checked by the kernel against the entry's data, and it was reviewed for fidelity before the proof attempt.
\end{abstract}
\maketitle

\section{Introduction}

For $n\ge0$ let
\[
  N_n = 3^{3^n}+2 .
\]
The OEIS entry A083020 \cite{A083020}, submitted by Cino Hilliard on May~31, 2003, is the irregular triangle whose row $n$ lists the prime factors of $N_n$ in increasing order.
Its example section reads $N_0=5$, $N_1=29$, $N_2=5\cdot31\cdot127$, $N_3=11\cdot693236134999$; the listed terms continue with the six prime factors of $N_4=3^{81}+2$, the $116$-digit number $N_5=3^{243}+2$ (so the entry asserts that $N_5$ is prime), and the two smallest prime factors $5,31$ of the $348$-digit number $N_6$.
The comments contain the conjecture that all of these numbers are squarefree (Conjecture~\ref{conj} below).

We show that the conjecture fails at $n=18$.

\begin{theorem}\label{thm:main}
$31^2=961$ divides $N_{18}=3^{3^{18}}+2$. In particular $N_{18}$ is not squarefree.
\end{theorem}

The number $N_{18}$ has $\lfloor 3^{18}\log_{10}3\rfloor+1=184{,}846{,}550$ decimal digits, so it is far out of reach of factorization; but divisibility by the fixed number $961$ only depends on the residue of $3^{3^{18}}$ modulo $961$, which Euler's theorem reduces to the residue of $3^{159}$ (Section~\ref{sec:proof}).
The same reduction describes all $n$ (Proposition~\ref{prop:all}): $961\mid N_n$ if and only if $n\equiv18\pmod{60}$, and $29^2=841$ divides $N_n$ if and only if $n\equiv73\pmod{84}$.
In particular infinitely many $N_n$ are not squarefree.

\subsection*{Status and prior work}
The entry A083020 (revision \#20 of January~4, 2026, which is still the current revision in October~2026) still states the conjecture; its EXTENSIONS section records only a correction by T.~D.~Noe (November~15, 2006) and a clarified name with the terms $a(14),a(15)$ added by Sean A.~Irvine (January~4, 2026).
The math0 literature referees found no prior counterexample or proof.
They checked the OEIS entries on related numbers, namely A055777 ($3^{3^n}$) \cite{A055777}, A129290 ($3^{3^n}+1$) \cite{A129290}, A129291 \cite{A129291}, A113913 \cite{A113913}, A140331 \cite{A140331} and A061678 \cite{A061678}, none of which discusses the squarefreeness of $N_n$; Stephan's collection \emph{Prove or disprove.\ 100 conjectures from the OEIS} \cite{St04}, which does not contain A083020; and web searches of arXiv, MathOverflow, Mathematics Stack Exchange and the SeqFan mailing list.
The conjecture is not among the OEIS statements of the formal-conjectures repository \cite{FC} at the commit pinned by math0 (\texttt{c77fdd4}).
We therefore classify Theorem~\ref{thm:main} as a new result, with the caveat that the argument is elementary and the counterexample may have been noticed without being recorded.
We do not claim that $n=18$ is the least $n$ for which $N_n$ is not squarefree (Remark~\ref{rem:small}).

\section{Statement}\label{sec:statement}

\subsection*{The OEIS conjecture}
The COMMENTS section of A083020 contains the following sentence, verbatim:
\begin{quote}
\raggedright{\small\texttt{Conjecture: These numbers have no repeating factors or 3\^{}(3\^{}n) + 2 is squarefree.}}
\end{quote}
The comment carries no separate attribution; it is part of the original entry by Cino Hilliard, May~31, 2003.
The entry has offset $0$, so ``these numbers'' are $N_n$ for all $n\ge0$, and the two phrasings joined by ``or'' say the same thing: $N_n$ has no repeated prime factor.

\begin{conjecture}[Hilliard]\label{conj}
For every integer $n\ge0$, the number $N_n=3^{3^n}+2$ is squarefree.
\end{conjecture}

By Theorem~\ref{thm:main}, Conjecture~\ref{conj} is false.

\subsection*{Lean statement}
The statement is the definition \decl{Math0.Statements.A083020.statement} in the math0 statement file \path{lean/Math0/Statements/OEIS/A083020.lean}, which was reviewed for fidelity and frozen (SHA-256 \texttt{d520c666\ldots ae473}) before any proof attempt.
The definitions that encode the triangle, and the statement, read verbatim (docstrings omitted; namespace \decl{Math0.Statements.A083020}):
\begin{lstlisting}
def N (n : ℕ) : ℕ := 3 ^ (3 ^ n) + 2

def row (n : ℕ) : List ℕ := Nat.primeFactorsList (N n)

def a (i : ℕ) : ℕ := ((List.range (i + 1)).flatMap row).getD i 0

def statement : Prop :=
  ∀ n : ℕ, Squarefree (3 ^ (3 ^ n) + 2)
\end{lstlisting}
In Mathlib, \texttt{Squarefree r} is defined as $\forall x,\ x\cdot x\mid r\to\texttt{IsUnit}\ x$, and the only unit of $\N$ is $1$ (\decl{Nat.isUnit_iff}); so for a natural number $r$ it says that no square $x^2$ with $x\ge2$ divides $r$, which is the usual notion.
Thus \decl{statement} is Conjecture~\ref{conj}, quantified over all $n\in\N$ as the offset $0$ requires, with no extra hypotheses; it involves no subtraction or division in $\N$, and $N_n\ge5$ rules out degenerate cases.
The file also proves
\begin{lstlisting}
theorem statement_iff : statement ↔ ∀ n, (row n).Nodup
\end{lstlisting}
from Mathlib's \decl{Nat.squarefree_iff_nodup_primeFactorsList}; this is the ``no repeating factors'' phrasing: no prime occurs twice in row $n$.
We prove \texttt{¬ statement}.

\subsection*{Encoding and check against the OEIS data}
Row $n$ of the triangle is \texttt{row n}, Mathlib's list of the prime factors of $N_n$ in increasing order with multiplicity. The listed rows contain no repeated primes, so listing distinct primes instead would give the same terms; the conjecture is stated directly about $N_n$ in any case.
The sequence itself is \texttt{a}, the triangle read by rows with offset $0$: \texttt{a i} is the entry at position $i$ of the concatenation of rows $0,\dots,i$, which has at least $i+1$ entries because every row is nonempty.

The function \texttt{a} is computable, but evaluating it means factoring $N_n$ by trial division, so the statement file checks the data through certificates instead.
The lemma \decl{primeFactorsList_eq_of} says that a sorted list of primes with product $m$ is \texttt{Nat.primeFactorsList m} (uniqueness of factorization), and \decl{row_data} applies it to
\[
  5,\quad 29,\quad 5\cdot31\cdot127,\quad 11\cdot 693236134999,\quad 5\cdot15679\cdot2372959\cdot5945531\cdot3312453269\cdot121032673159
\]
for $n=0,\dots,4$. The products are checked by \texttt{norm\_num}, sortedness by \texttt{decide}, primes below $10^5$ by \texttt{norm\_num}, and the larger primes by Lucas certificates (Table~\ref{tab:lucas}).
Mathlib's \decl{lucas_primality} says that $p$ is prime if some $w$ satisfies $w^{p-1}\equiv1\pmod p$ and $w^{(p-1)/q}\not\equiv1\pmod p$ for every prime $q\mid p-1$.
The Boolean function \decl{lucasCheck} tests these conditions for a given $w$ and a list of primes with product $p-1$, computing powers with a square-and-multiply function \decl{powMod} that is proved equal to $w^e\bmod p$ for $e<2^{64}$; the theorem \decl{prime_of_lucasCheck} turns a successful check into a proof of primality, and each check is evaluated by \texttt{decide +kernel}.
The primes in the lists for $p-1$ are proved prime by \texttt{norm\_num} or by certificates of their own.
The theorem \decl{a_data} then states that $a(0),\dots,a(12)$ are the $13$ listed terms $5, 29, 5, 31, 127, 11, 693236134999, 5, 15679, 2372959, 5945531, 3312453269, 121032673159$.

\begin{table}[h]
\centering\small
\begin{tabular}{rrl}
\toprule
$p$ & $w$ & factorization of $p-1$\\
\midrule
$391151$ & $13$ & $2\cdot5^2\cdot7823$\\
$693236134999$ & $7$ & $2\cdot3^2\cdot11\cdot8951\cdot391151$\\
$2372959$ & $3$ & $2\cdot3^2\cdot7\cdot37\cdot509$\\
$5945531$ & $2$ & $2\cdot5\cdot37\cdot16069$\\
$22381441$ & $7$ & $2^7\cdot3\cdot5\cdot11657$\\
$3312453269$ & $2$ & $2^2\cdot37\cdot22381441$\\
$4317661$ & $6$ & $2^2\cdot3^2\cdot5\cdot17^2\cdot83$\\
$20172112193$ & $3$ & $2^6\cdot73\cdot4317661$\\
$121032673159$ & $3$ & $2\cdot3\cdot20172112193$\\
\bottomrule
\end{tabular}
\caption{The Lucas certificates in the statement file: for each $p$, the witness $w$ and the prime factors of $p-1$.}
\label{tab:lucas}
\end{table}

The remaining listed terms are $a(13)$, a $116$-digit number, and $a(14)=5$, $a(15)=31$.
The theorem \decl{N5_data} checks that $a(13)$ as listed equals $N_5=3^{243}+2$, and \decl{row6_data} checks that row $6$ begins with $5,31$, by computing the two smallest prime factors of $N_6$.
The primality of $N_5$ is not proved in Lean (SymPy's primality test reports it prime), so $a(13),a(14),a(15)$ sit at the listed positions exactly when $N_5$ is prime, and these positions are not checked as indices.
All of these checks live in the statement file; the refutation below does not depend on them.

\section{Proof}\label{sec:proof}

The proof follows the declarations of the Lean file: the private lemmas \decl{pow_930}, \decl{split_exp} and \decl{sq_dvd}, and the theorem \decl{result}.
Congruences are in $\N$ (\texttt{Nat.ModEq}).

\begin{lemma}[\decl{pow_930}]\label{lem:euler}
$3^{930}\equiv1\pmod{961}$.
\end{lemma}

\begin{proof}
Since $31$ is prime, $\varphi(31^2)=31\cdot30=930$, and $\gcd(3,961)=1$, so this is Euler's theorem.
The Lean proof does not use Euler's theorem: it unfolds the congruence to $3^{930}\bmod 961=1\bmod 961$ and checks this equation by \texttt{decide +kernel}, which makes the kernel evaluate both sides; $3^{930}$ has $1475$ binary digits.
\end{proof}

\begin{lemma}[\decl{split_exp}]\label{lem:split}
$3^{18}=930\cdot416581+159$, and hence $3^{3^{18}}=\bigl(3^{930}\bigr)^{416581}\cdot3^{159}$.
\end{lemma}

\begin{proof}
$3^{18}=387420489$ and $930\cdot416581=387420330$.
The second claim follows from the first by the laws $(x^b)^c=x^{bc}$ and $x^bx^c=x^{b+c}$.
In Lean, the right-hand side is rewritten with these two laws (\decl{pow_mul}, \decl{pow_add}) into $3^{930\cdot416581+159}$, and \texttt{norm\_num} checks that the two exponents agree; the number $3^{3^{18}}$ is never evaluated.
\end{proof}

\begin{lemma}[\decl{sq_dvd}]\label{lem:dvd}
$31\cdot31$ divides $3^{3^{18}}+2$.
\end{lemma}

\begin{proof}
Raising Lemma~\ref{lem:euler} to the power $416581$, multiplying by $3^{159}$ on the right and adding $2$ gives
\[
  \bigl(3^{930}\bigr)^{416581}\cdot3^{159}+2\equiv 1^{416581}\cdot3^{159}+2\pmod{961}
\]
(\decl{Nat.ModEq.pow}, \decl{Nat.ModEq.mul_right}, \decl{Nat.ModEq.add_right}).
Next, $3^{159}\equiv959\equiv-2\pmod{961}$, so $1^{416581}\cdot 3^{159}+2\equiv961\equiv0\pmod{961}$.
To check the residue by hand: the successive squares $3^{2^k}\bmod961$ for $k=0,\dots,7$ are
\[
  3,\ 9,\ 81,\ 795,\ 648,\ 908,\ 887,\ 671,
\]
and $159=2^7+2^4+2^3+2^2+2+1$, so $3^{159}\equiv3\cdot9\cdot81\cdot795\cdot648\cdot671$; the partial products reduce to $27,\ 265,\ 216,\ 623,\ 959$ modulo $961$.
In Lean, the congruence $1^{416581}\cdot3^{159}+2\equiv0\pmod{961}$ is checked by \texttt{decide +kernel}.
By transitivity and Lemma~\ref{lem:split}, $3^{3^{18}}+2\equiv0\pmod{961}$, that is, $961$ divides $3^{3^{18}}+2$ (\decl{Nat.modEq_zero_iff_dvd}).
In Lean the goal $31\cdot31\mid\dots$ is first rewritten with Lemma~\ref{lem:split}, and $31\cdot31$ is identified with $961$ by evaluating the product of the two literals.
\end{proof}

As a consistency check, which the proof does not use: modulo $31$ we have $3^{30}\equiv1$ and $3^{18}\equiv9\pmod{30}$, so $3^{3^{18}}\equiv3^9\equiv29\equiv-2\pmod{31}$, as Lemma~\ref{lem:dvd} requires.

\begin{proof}[Proof of Theorem~\ref{thm:main} and of \texttt{¬ statement} (\decl{result})]
Lemma~\ref{lem:dvd} is the first claim.
Suppose that \decl{statement} holds. Then $N_{18}$ is squarefree, so by the definition of \texttt{Squarefree}, applied to $x=31$ and Lemma~\ref{lem:dvd}, $31$ is a unit of $\N$.
By \decl{Nat.isUnit_iff} this means $31=1$, which is false (\texttt{omega}).
\end{proof}

\subsection*{All \texorpdfstring{$n$}{n}, and small primes}
The following is not part of the formal proof.

\begin{proposition}\label{prop:all}
For $n\ge0$: $961\mid N_n$ if and only if $n\equiv18\pmod{60}$, and $841\mid N_n$ if and only if $n\equiv73\pmod{84}$.
In particular infinitely many $N_n$ are not squarefree.
\end{proposition}

\begin{proof}
Let $x_n=3^{3^n}\bmod961$, so that $x_{n+1}\equiv x_n^3\pmod{961}$, and $961\mid N_n$ if and only if $x_n=959$.
By Lemma~\ref{lem:euler}, $x_n$ depends only on $3^n\bmod930$.
Now $930=3\cdot310$ with $\gcd(3,310)=1$. For $n\ge1$ we have $3^n\equiv0\pmod 3$, and $3^n\bmod310$ is periodic in $n$ with period the order of $3$ modulo $310=2\cdot5\cdot31$, which is $\operatorname{lcm}(1,4,30)=60$.
By the Chinese remainder theorem, $3^n\bmod930$, and hence $x_n$, depends only on $n\bmod60$ for $n\ge1$.
Computing $x_1,\dots,x_{60}$ by repeated cubing modulo $961$ gives $x_n=959$ only for $n=18$, and $x_0=3$.
This proves the first claim.
For $841=29^2$ we have $\varphi(841)=812=2^2\cdot7\cdot29$, which is prime to $3$, so $3^n\bmod812$ is purely periodic in $n$ with period the order of $3$ modulo $812$, which is $\operatorname{lcm}(2,6,28)=84$. Hence $y_n=3^{3^n}\bmod841$ is periodic with period dividing $84$, and computing $y_0,\dots,y_{83}$ gives $y_n=839$ only for $n=73$.
\end{proof}

\begin{remark}[Smaller $n$]\label{rem:small}
For $n\le4$ the factorizations in \decl{row_data} show that $N_n$ is squarefree, and $N_5$ is squarefree if it is prime, as the entry asserts.
For $6\le n\le17$ we do not know whether $N_n$ is squarefree: these numbers have between $348$ and about $6.2\times10^7$ decimal digits, and none of them is completely factored (for $N_6$ only the factors $5$ and $31$ and a composite cofactor are known).
A computer search (not part of the formal proof) that computes $3^{3^n}\bmod p^2$ by repeated cubing, for every prime $p<10^8$ and every $n\le18$, finds $p^2\mid N_n$ only for $(n,p)=(18,31)$.
So $n=18$ is the least $n$ for which $N_n$ is divisible by the square of a prime below $10^8$, but we do not know whether it is the least $n$ for which $N_n$ is not squarefree.
\end{remark}

\section{Machine verification}\label{sec:lean}

\subsection*{Files and declarations}
The proof is in the file \path{lean/Math0/Frontier/OEIS/A083020.lean} of the math0 repository, \url{https://github.com/Foundation-0/math0}.
It imports the frozen statement file \path{lean/Math0/Statements/OEIS/A083020.lean} (Section~\ref{sec:statement}) and the audit commands in \path{lean/Math0/Audit.lean}.
The main declaration is \decl{Math0.Frontier.A083020.result}, of type \texttt{¬ Math0.Statements.A083020.statement}; the auxiliary lemmas \decl{pow_930}, \decl{split_exp} and \decl{sq_dvd} are private.
The proof file reads, verbatim except for the license header, the imports, the docstrings and blank lines:
\begin{lstlisting}
namespace Math0.Frontier.A083020
open Math0.Statements.A083020
private theorem pow_930 : 3 ^ 930 ≡ 1 [MOD 961] := by
  unfold Nat.ModEq
  decide +kernel
private theorem split_exp : (3 : ℕ) ^ (3 ^ 18) = (3 ^ 930) ^ 416581 * 3 ^ 159 := by
  rw [← pow_mul, ← pow_add]
  norm_num
private theorem sq_dvd : 31 * 31 ∣ 3 ^ (3 ^ 18) + 2 := by
  have h : (3 ^ 930) ^ 416581 * 3 ^ 159 + 2 ≡ 1 ^ 416581 * 3 ^ 159 + 2 [MOD 961] :=
    ((pow_930.pow 416581).mul_right _).add_right _
  have h0 : 1 ^ 416581 * 3 ^ 159 + 2 ≡ 0 [MOD 961] := by
    unfold Nat.ModEq
    decide +kernel
  rw [split_exp]
  exact Nat.modEq_zero_iff_dvd.1 (h.trans h0)
theorem result : ¬ statement := by
  intro h
  have hu : IsUnit (31 : ℕ) := h 18 31 sq_dvd
  rw [Nat.isUnit_iff] at hu
  omega
end Math0.Frontier.A083020
#math0_settles Math0.Frontier.A083020.result Math0.Statements.A083020.statement
\end{lstlisting}
The file uses no \texttt{sorry}, \texttt{native\_decide}, \texttt{set\_option} or macros.

\subsection*{How the kernel checks the computations}
The tactic \texttt{decide +kernel} proves a decidable proposition by asking the kernel to reduce its decision procedure to \texttt{true}.
Lean's kernel evaluates arithmetic on natural number literals (addition, multiplication, powers, division and remainder, among others) with arbitrary-precision integer arithmetic, so the checks $3^{930}\bmod961=1$ and $(1^{416581}\cdot3^{159}+2)\bmod961=0$ are immediate; unlike \texttt{native\_decide}, this adds no axiom.
The only large exponent is handled symbolically: the proof manipulates \texttt{Nat.ModEq} terms whose sides are never evaluated, and Lemma~\ref{lem:split} only compares the exponents $3^{18}$ and $930\cdot416581+159$.
The data checks in the statement file (Section~\ref{sec:statement}) use \texttt{decide +kernel} in the same way, on the Lucas certificates.

\subsection*{Audit}
The command \texttt{\#math0\_settles}, defined in \path{lean/Math0/Audit.lean} for statements that math0 formalizes itself, fails unless (1) \decl{statement} is a \texttt{Prop}-valued definition without universe parameters; (2) both \decl{statement} and \decl{result} use no axioms besides \decl{propext}, \decl{Classical.choice} and \decl{Quot.sound}, which rejects \decl{sorryAx}, the axiom \decl{Lean.ofReduceBool} introduced by \texttt{native\_decide}, and custom axioms; and (3) the type of \decl{result} is syntactically the statement or its negation, up to metadata annotations, with no unfolding.
Running \texttt{lake env lean Math0/Frontier/OEIS/A083020.lean} in the directory \path{lean} prints\par\nobreak\noindent
\begin{minipage}{\linewidth}
\begin{lstlisting}[basicstyle=\ttfamily\scriptsize]
math0 audit: 'Math0.Statements.A083020.statement' depends on axioms [propext]
math0 audit: 'Math0.Frontier.A083020.result' depends on axioms [Quot.sound, propext]
math0 statement check: 'Math0.Frontier.A083020.result' refutes 'Math0.Statements.A083020.statement'
\end{lstlisting}
\end{minipage}
so \decl{result} proves the negation of the statement and depends only on \decl{propext} and \decl{Quot.sound}, a subset of the three permitted axioms.

\subsection*{Toolchain and reproduction}
We use Lean~4 \cite{Lean4} \texttt{v4.33.1} and Mathlib \cite{Mathlib} \texttt{v4.33.1}, pinned in \path{lean/}; the repository also pins formal-conjectures \cite{FC} at commit \texttt{c77fdd4}, which this result does not use.
To reproduce, run \path{scripts/verify.sh} from the repository root: it fetches the Mathlib build cache, builds all math0 results with \texttt{lake build} (re-running every audit command), checks that the pinned \texttt{formal\_conjectures} and \texttt{mathlib} checkouts are unmodified, and reports the number of passed statement checks.
To check this result alone, run \texttt{lake build Math0.Statements.OEIS.A083020} and then \texttt{lake env lean Math0/Frontier/OEIS/A083020.lean} in the directory \path{lean}; the kernel computations take a few seconds, and most of the running time is spent loading Mathlib.

{\small
\begin{thebibliography}{99}

\bibitem{FC}
M.~Firsching, P.~Lezeau, S.~Mercuri, M.~Z. Horv\'ath, Y.~Dillies, C.~S\"onne, E.~Wieser, F.~Zhang, T.~Hubert, B.~Ag\"uera y Arcas, and P.~Kohli, \emph{Formal Conjectures: An Open and Evolving Benchmark for Verified Discovery in Mathematics}, arXiv:2605.13171, 2026. Repository: \url{https://github.com/google-deepmind/formal-conjectures}.

\bibitem{Mathlib}
The mathlib Community, \emph{The Lean mathematical library}, in: Proceedings of the 9th ACM SIGPLAN International Conference on Certified Programs and Proofs (CPP 2020), ACM, 2020, 367--381.

\bibitem{Lean4}
L.~de Moura and S.~Ullrich, \emph{The Lean 4 theorem prover and programming language}, in: Automated Deduction -- CADE 28, Lecture Notes in Comput. Sci. \textbf{12699}, Springer, 2021, 625--635.

\bibitem{A055777}
The OEIS Foundation Inc., \url{https://oeis.org/A055777} ($3^{3^n}$), entry by H.~Bottomley, Jul.~12, 2000.

\bibitem{A061678}
The OEIS Foundation Inc., \url{https://oeis.org/A061678} (continued fraction for $\sum_{k\ge0}3^{-3^k}$), entry by J.~Earls, Jun.~23, 2001.

\bibitem{A083020}
The OEIS Foundation Inc., \url{https://oeis.org/A083020} (prime factors of $3^{3^n}+2$), entry by C.~Hilliard, May~31, 2003; revision \#20, Jan.~4, 2026.

\bibitem{A113913}
The OEIS Foundation Inc., \url{https://oeis.org/A113913} (number of distinct prime factors of $3^{3^n}+1$), entry by J.~Coveiro, Jan.~29, 2006.

\bibitem{A129290}
The OEIS Foundation Inc., \url{https://oeis.org/A129290} (base-$3$ Fermat numbers $3^{3^n}+1$), entry by A.~Adamchuk, Apr.~8, 2007.

\bibitem{A129291}
The OEIS Foundation Inc., \url{https://oeis.org/A129291} ($1-3^{3^n}+9^{3^n}$), entry by A.~Adamchuk, Apr.~8, 2007.

\bibitem{A140331}
The OEIS Foundation Inc., \url{https://oeis.org/A140331} (smallest $k$ such that $3^{3^n}-k$ is prime), entry by J.~Vos Post, May~26, 2008.

\bibitem{St04}
R.~Stephan, \emph{Prove or disprove. 100 conjectures from the OEIS}, arXiv:math/0409509, 2004.

\end{thebibliography}
}

\end{document}
